一个算法找到目标值 \(12.3\) 的可行解，只能说明最优值不高于 \(12.3\)；也许还有目标值 \(5\) 的方案尚未发现。要证明候选解接近最优，还需要从另一侧给出下界。Lagrange 对偶的核心正是系统地制造这种下界。

本单元沿``下界---间隙---最优性方程---现实解释''推进：

\begin{enumerate}
\tightlist
\item
  Lagrangian 怎样制造全局下界从约束违规项的符号出发推导对偶函数和弱对偶。
\item
  对偶间隙何时闭合区分弱对偶与强对偶，并解释 Slater 条件为何排除病态边界。
\item
  KKT 把最优性写成一组方程把原始可行性、对偶可行性、驻点与互补松弛组合成可验证证书。
\item
  对偶变量是约束的影子价格通过扰动价值函数说明乘子的单位与现实含义，并把原始---对偶间隙用于算法停止。
\end{enumerate}

符号在不同教材中可能因约束方向而改变。本单元统一使用最小化问题和 \(f_i(x)\le0\)；在这个约定下，不等式乘子必须非负。若改用 \(\ge0\)，相应符号也必须一起改变。

\subsection{一个简单的对偶练习}

考虑 \(\min\{x^2:x\ge1\}\)。原始解 \(x^\star=1,\;p^\star=1\)。将约束写成 \(1-x\le0\)，乘子 \(\lambda\ge0\)，Lagrangian 为 \(L(x,\lambda)=x^2+\lambda(1-x)\)。对 \(x\) 最小化得 \(g(\lambda)=\lambda-\lambda^2/4\)。最大化 \(g(\lambda)\) 得 \(\lambda^\star=2,\;d^\star=1\)。于是 \(d^\star=p^\star=1\)，上下界闭合。这个数例包含了对偶理论的全部关键步骤：构造下界、弱对偶自动成立、在凸条件和严格可行下强对偶闭合。

\subsection{怎么读这一章}

先把对偶看成一种证明办法：原问题在寻找尽可能小的可行值，对偶问题则努力制造一个任何可行解都不可能低于的下界。第一篇说明下界从哪里来，第二篇讨论上下界何时相遇，第三篇把相遇条件写成 KKT，最后一篇解释乘子为何能衡量放宽约束的价值。标准问题形式可回看 \hyperref[ch:074]{``识别凸优化问题''}。

\subsection{本章篇目}

以下篇目按概念依赖排列。

\begin{itemize}
\tightlist
\item \hyperref[sec:09-Convex-Optimization/05-Duality/01]{``Lagrangian 怎样制造全局下界''}；
\item \hyperref[sec:09-Convex-Optimization/05-Duality/02]{``对偶间隙何时闭合''}；
\item \hyperref[sec:09-Convex-Optimization/05-Duality/03]{``KKT 把最优性写成一组方程''}；
\item \hyperref[sec:09-Convex-Optimization/05-Duality/04]{``对偶变量是约束的影子价格''}。
\end{itemize}
